% !PRINTED 550
\weberpage{550}
%<<CONT
\weq{17}{\sqrt[m]{\psi_{\lambda}} = \chi_{\lambda}\left(\sqrt{a_1^{m_1} a_2^{m_2} \ldots}\right)^{-\nu} e^{-\frac{i\Theta}{m}\nu}.}

Nach (15) können wir setzen
\[
e^{i\Theta_1} = \frac{b_1 + ic_1}{\sqrt{a_1^m}}, \qquad e^{i\Theta_2} = \frac{b_2 + ic_2}{\sqrt{a_2^m}} \cdots,
\]
und wenn wir also durch Zerlegung in den reellen und imaginären Bestandtheil
\weq{18}{(b_1 + ic_1)^{m_1}(b_2 + ic_2)^{m_2} \cdots = B + iC}
erhalten und
\weq{19}{a_1^{m_1} a_2^{m_2} \cdots = A}
setzen, so folgt nach (16)
\[
e^{i\Theta} = \frac{B + iC}{\sqrt{A^m}},
\]
also
\weq{20}{\sqrt{A^m}\cos\Theta = B, \qquad \sqrt{A^m}\sin\Theta = C,}
und
\weq{21}{\sqrt[m]{\psi_{\lambda}} = \chi_{\lambda}\sqrt{A}^{-\nu}\left(\cos\frac{\Theta\nu}{m} - i\sin\frac{\Theta\nu}{m}\right).}

Der Winkel $\Theta$ ist durch (20) nur bis auf ein Vielfaches von
$2\pi$ bestimmt, und wenn man also $\Theta + 2h\pi$ für $\Theta$ setzt und $h$
von $0$ bis $m - 1$ gehen lässt, so erhält man aus (21) die $m$ verschiedenen Werthe der Wurzelgrösse. Die Functionen $\cos\dfrac{\Theta\nu}{m}$
und $\sin\dfrac{\Theta\nu}{m}$ können noch rational durch $\cos\dfrac{\Theta}{m}$, $\sin\dfrac{\Theta}{m}$ ausgedrückt werden. Die Auflösung der cyklischen Gleichungen in
dem reellen Körper $\Omega$ ist also auf Folgendes zurückgeführt.

\gesperrt{Man adjungirt zunächst dem Körper $\Omega$ die $m^{\mathrm{te}}$ Einheitswurzel $\varepsilon$ (Theilung der Kreisperipherie in $m$ gleiche
Theile). Hierauf sind $A, B, C$ bekannt. Man adjungirt
ferner die positive Quadratwurzel $\sqrt{A}$, dann sind
$\cos\Theta$ und $\sin\Theta$ durch (20) bekannt. Endlich adjungirt
man $\cos\dfrac{\Theta}{m}$, $\sin\dfrac{\Theta}{m}$ (Theilung des Winkels $\Theta$ in $m$ Theile).
Dann ist die cyklische Gleichung durch (21) gelöst.}

Diese Betrachtungen führen noch zu einem interessanten
Resultat über die Realitätsverhältnisse der Wurzeln cyklischer
Gleichungen.
